\subsection{Strategie promptování}
\label{subsec:prompts-strategies}

Způsob, jakým je prompt formulován, zásadně ovlivňuje kvalitu a konzistentnost výstupu modelu.
V literatuře jsou popsány tři základní strategie promptování, které se liší mírou kontextu a strukturou poskytnutou modelu~\cite{brown-gpt3,wei-cot}.

\subsubsection*{Zero-shot prompting}
\label{subsubsec:prompts-strategies-zero-shot}

Při zero-shot přístupu je modelu zadán úkol bez jakéhokoli příkladu očekávaného výstupu~\cite{kojima-zeroshot}.
Model se spoléhá výhradně na znalosti získané při trénování.
Pro analýzu studentského kódu by zero-shot prompt vypadal přibližně takto: \uv{Analyzuj následující zdrojový kód a identifikuj v něm problémy.}

Výhodou tohoto přístupu je jednoduchost a nízká spotřeba tokenů.
Nevýhodou je, že model nemá explicitní vodítko pro formát a úroveň detailu odpovědi, což může vést k nekonzistentním výstupům.
Jeden dotaz může vrátit stručný seznam problémů, zatímco jiný rozsáhlé vysvětlení s příklady oprav.

\subsubsection*{Few-shot prompting}
\label{subsubsec:prompts-strategies-few-shot}

Few-shot přístup rozšiřuje prompt o jeden nebo více příkladů (vzorových dvojic vstup/výstup), které ukazují modelu požadovaný formát a úroveň detailu odpovědi~\cite{brown-gpt3}.
Model tak může na základě poskytnutých příkladů lépe pochopit, jaký typ výstupu je očekáván.

Pro systém Kelvin by few-shot prompt mohl obsahovat ukázkový zdrojový kód s předdefinovanými komentáři, které ukazují, jakým způsobem má model identifikovat a popsat nalezené problémy.
Nevýhodou je zvýšená spotřeba tokenů, protože každý příklad zabírá prostor v kontextovém okně.
U modelů s menším kontextovým oknem může tento přístup významně omezit prostor pro samotný analyzovaný kód.

\subsubsection*{Chain-of-thought prompting}
\label{subsubsec:prompts-strategies-chain-of-thought}

Strategie CoT vede model k tomu, aby svou analýzu provedl v několika po sobě jdoucích krocích~\cite{wei-cot}.
Místo přímého vygenerování konečného výstupu model nejprve provede hrubou analýzu, poté ji kriticky zhodnotí a na základě tohoto zhodnocení vytvoří finální seznam zjištěných problémů.
Tato logika lze jednoduše nahradit za zero-shot a to v případě, že model poskytuje \uv{thinking} logiku, která pracuje obdobně jako CoT\@.

V kontextu systému Kelvin byl navržen třístupňový CoT přístup:

\begin{enumerate}
    \item \textbf{Hrubá analýza} (\textit{draft analysis}) identifikuje všechny potenciální problémy v kódu studenta.
    V tomto kroku je model instruován, aby byl co nejdůkladnější a nezabýval se filtrováním irelevantních nálezů.

    \item \textbf{Kritická revize} (\textit{critique analysis}) zpracuje výstup předchozího kroku a vyhodnotí relevanci každého nálezu vzhledem k zadání úlohy.
    Tento krok slouží k odstranění falešných pozitiv a nálezů, které jsou sice technicky správné, ale v kontextu zadání nejsou důležité.

    \item \textbf{Finální přehled} (\textit{review analysis}) vytvoří konečný strukturovaný výstup obsahující pouze relevantní a ověřené problémy ve formátu vhodném pro systém.
\end{enumerate}

CoT přístup je výpočetně náročnější, protože vyžaduje tři po sobě jdoucí volání modelu.
Jak bylo ukázáno v \tabref{tab:cloud-api-costs}, náklady na jedno odevzdání při CoT přístupu jsou přibližně 4 až 10krát vyšší než u přímého přístupu s jedním voláním.
Zároveň je zpracování jednoho odevzdání výrazně pomalejší.
Na druhou stranu CoT přístup může vést k přesnějším a lépe strukturovaným výstupům, protože každý krok slouží jako kontrolní mechanismus pro výstup předchozího kroku.

\endinput